\documentclass[11pt]{amsart}
\usepackage{amsmath,amssymb}
\usepackage[hidelinks]{hyperref}

\title{Literature answer for the dual ASQ versus UASQ question}
\author{fa\_banach\_001, agent\_lane\_04}
\date{2026-06-29}

\begin{document}
\maketitle

\section*{Status}

This is a literature-status packet, not a new proof.  It records that the
first open question in Garcia-Lirola--Rueda Zoca \cite{GLRZ2017} is answered
by the later paper of Abrahamsen--Hajek--Troyanski \cite{AHT2019}.

\section*{Source question}

In Section 6, Question 6.1 of \cite{GLRZ2017}, labelled
\texttt{UASQASQequivalent} in the arXiv source, the authors ask:
\begin{quote}
Are UASQ and ASQ equivalent, at least, for dual Banach spaces?
\end{quote}
The same paper proves that no dual Banach space is UASQ.  Thus any dual ASQ
Banach space gives a negative answer.

\section*{Supporting answer}

The introduction of \cite{AHT2019} explicitly discusses this question: after
recalling UASQ, it says that Question 6.1 in the ASQ/UASQ literature asked
whether dual ASQ spaces are UASQ, and that the construction of a dual ASQ
space answers this in the negative.  The decisive theorem is
\cite[Theorem 3.7]{AHT2019}: a dual Banach space admits a dual ASQ norm if and
only if it contains an isomorphic copy of $c_0$.

Combining that theorem with the no-dual-UASQ theorem from \cite{GLRZ2017}
gives the answer to the source question:
\[
        \text{dual ASQ} \not\Rightarrow \text{UASQ}.
\]
Therefore UASQ and ASQ are not equivalent for dual Banach spaces.

\section*{Scope}

This packet records only the first Section 6 question of \cite{GLRZ2017}.  It
does not answer the later questions in the same source about uniformly
discrete metric spaces $M$ with $Lip(M)$ ASQ, or about isomorphic duality of
$S(M,X)$.

\begin{thebibliography}{2}

\bibitem{GLRZ2017}
L. Garcia-Lirola and A. Rueda Zoca,
\emph{Unconditional almost squareness and applications to spaces of Lipschitz
functions},
J. Math. Anal. Appl. 451 (2017), 117--131; arXiv:1605.04699.

\bibitem{AHT2019}
T. A. Abrahamsen, P. Hajek, and S. Troyanski,
\emph{Almost square dual Banach spaces},
arXiv:1912.12519.

\end{thebibliography}

\end{document}
